\section{重参数化让随机节点可以反向传播}
\label{sec:14-Deep-Learning/06-Variational-Autoencoders/02}

上一节把 ELBO 写了出来，也顺手交代了重参数化的公式。现在把那一步单独拆开看，因为它藏着一个很容易被跳过的困难。

假设你已经照着 ELBO 搭好了计算图：encoder 吃进 \(x\)，输出 \(\mu_\phi(x)\) 和 \(\log\sigma_\phi^2(x)\)；从这个高斯里采一个 \(z\)；decoder 吃进 \(z\)，输出重构似然。现在调用一次反向传播。decoder 的参数 \(\theta\) 拿到了梯度，encoder 的参数 \(\phi\) 拿到的是零——或者更常见的情形，框架直接报错说采样这个算子没有定义导数。

问题不在实现。写出目标的那一刻它就在了：

\[
\mathcal L(\phi) = \E_{q_\phi(z\given x)}[f(z)] .
\]

这里的 \(f\) 是 ELBO 的被积部分，也可以是任何依赖 \(z\) 的标量分数，具体形式暂时不重要。要紧的是参数 \(\phi\) 出现在\emph{期望的下标里}，不是在被积函数里。展开成积分就清楚了：

\[
\mathcal L(\phi) = \int q_\phi(z\given x) f(z)\,dz.
\]

\(\phi\) 变一点，改变的是每个 \(z\) 被抽中的权重，而不是 \(f\) 在某个固定 \(z\) 上的取值。可你手里只有若干个已经抽出来的样本 \(z^{(1)},\dots,z^{(m)}\)——从这些样本上，你看不见``如果 \(\phi\) 再大一点，这个 \(z\) 本来会更容易被抽中''。采样这一步把 \(\phi\) 的信息吃掉了。

绕过它有两条路，走向完全不同。

\subsection{路线一：把参数搬回被积函数}

第一条路只用一个恒等式。对 \(q_\phi\) 求导，再乘除一个 \(q_\phi\)：

\[
\begin{aligned}
\nabla_\phi \int q_\phi(z) f(z)\,dz
&= \int \nabla_\phi q_\phi(z)\, f(z)\,dz\\
&= \int q_\phi(z)\,\frac{\nabla_\phi q_\phi(z)}{q_\phi(z)}\, f(z)\,dz\\
&= \E_{q_\phi}\!\left[f(z)\,\nabla_\phi \log q_\phi(z)\right].
\end{aligned}
\]

右端又是一个对 \(q_\phi\) 的期望，可以直接用样本平均估计。这就是得分函数估计量，强化学习里叫 REINFORCE，\(\nabla_\phi\log q_\phi\) 是统计里的得分函数。

它的适用范围宽得惊人：整个推导里 \(f\) 只是被乘上去的一个数，从来没有对 \(f\) 求过导。\(f\) 可以不连续，可以是黑箱模拟器的返回值，\(z\) 可以是离散的类别标签——只要 \(q_\phi\) 对 \(\phi\) 可微、密度写得出来就行。VAE 里想用离散隐变量（比如离散的笔画类型），这几乎是唯一能直接上手的办法。

问题写在同一个式子里。梯度的每个样本项是 \(f(z^{(i)})\nabla_\phi\log q_\phi(z^{(i)})\)：\(f\) 的\emph{整体尺度}直接乘在方向上。把 \(f\) 平移一个常数 \(c\)，真实梯度纹丝不动（因为 \(\E_q[\nabla_\phi\log q_\phi]=0\)），但估计量的方差会随 \(c^2\) 增长。重构对数似然动辄是几百量级的负数，这个常数偏移带来的噪声可以把有用信号淹没几个数量级。

标准补救是扣一个基线：用 \(f(z)-b\) 代替 \(f(z)\)。只要 \(b\) 与 \(z\) 无关，由同一个恒等式 \(\E_q[b\,\nabla_\phi\log q_\phi]=0\)，估计量仍然无偏，而方差可以大幅下降。基线通常取 \(f\) 的滑动平均，或者训练一个小网络去预测它。这是让得分函数估计在实践中可用的必要工序，不是可选优化。

\subsection{路线二：把随机性搬出参数}

第二条路换了个角度：既然麻烦在于``分布依赖参数''，那就想办法让分布不依赖参数。

对角高斯提供了一个干净的分解。从 \(\mathcal N(\mu,\sigma^2)\) 采样，等价于从与参数无关的 \(\mathcal N(0,1)\) 采样再做一次仿射变换：

\[
\varepsilon \sim \mathcal N(0,I),
\qquad
z = g_\phi(x,\varepsilon) = \mu_\phi(x) + \sigma_\phi(x)\odot\varepsilon .
\]

采样发生在 \(\varepsilon\) 上，\(\varepsilon\) 的分布里没有 \(\phi\)。于是目标可以重写成对一个\emph{固定}分布的期望：

\[
\mathcal L(\phi) = \E_{\varepsilon\sim \mathcal N(0,I)}\big[f\big(g_\phi(x,\varepsilon)\big)\big],
\]

这一次，参数只出现在被积函数里。梯度算子和期望换序之后，链式法则可以一路穿下去：

\[
\nabla_\phi \mathcal L(\phi)
= \E_{\varepsilon}\!\left[\nabla_z f(z)\big|_{z=g_\phi(x,\varepsilon)}\cdot \nabla_\phi g_\phi(x,\varepsilon)\right].
\]

计算图上的表现是：\(\varepsilon\) 成了一个外部输入的常数张量，从 \(\varepsilon\) 到 \(z\) 的那条边是普通的乘加运算，反向传播照常通过。原来那个``不可导的采样节点''消失了。

方差为什么小得多，从两个式子的形状就能看出来。得分函数估计量里出现的是 \(f\) 的\emph{取值}，重参数化估计量里出现的是 \(\nabla_z f\)，也就是 \(f\) 的\emph{斜率}。前者对 \(f\) 加常数敏感，后者完全免疫；前者只知道``这次抽到的 \(z\) 得分是 300''，后者知道``往哪个方向挪一点 \(z\) 会让得分升高''。后一种信息在单个样本上就有方向性，所以 VAE 训练时每步抽一个 \(\varepsilon\) 通常就够用。

\begin{center}
\begin{tikzpicture}[x=1.05cm,y=1cm,>=stealth]
  \draw[->] (-5.4,0) -- (5.6,0) node[right,font=\small] {梯度取值};
  \draw[->] (-5.4,0) -- (-5.4,5.4) node[above,font=\small] {频次};
  \foreach \tk in {-4,-2,0,2,4}{\draw (\tk,0) -- (\tk,-0.12) node[below,font=\scriptsize] {\(\tk\)};}
  \draw[thick,blue!60!black,fill=blue!12,fill opacity=0.55]
    plot[domain=-5:5,samples=140] (\x,{1.1*exp(-(\x-1)^2/9.68)}) -- (5,0) -- (-5,0) -- cycle;
  \draw[thick,red!65!black,fill=red!16,fill opacity=0.65]
    plot[domain=-0.8:2.8,samples=120] (\x,{4.4*exp(-(\x-1)^2/0.605)}) -- (2.8,0) -- (-0.8,0) -- cycle;
  \draw[dashed,black!55] (1,0) -- (1,5.05);
  \node[above,font=\scriptsize] at (1,5.05) {真实梯度 \(\nabla_\phi\mathcal L\)};
  \node[font=\small,blue!45!black] at (-3.3,1.35) {得分函数估计};
  \node[font=\small,red!55!black,anchor=west] at (2.0,3.6) {重参数化};
\end{tikzpicture}

\emph{同一个目标、同一个真实梯度，两种估计量的抽样分布宽窄相差一个量级；无偏说的是两条钟形曲线共用一个均值，方差说的是单次抽样离这个均值有多远。}
\end{center}

图里两条曲线均值重合，这一点值得停一下：两种估计都是无偏的，所以``哪个更准''不能靠期望回答。差别全在宽度上，而宽度决定了你需要多少样本、能用多大的学习率、以及训练会不会看起来像在原地打转。

\subsection{换序凭什么合法}

上面那步``梯度算子和期望换序''写起来只有一行，它是整条路线的承重点。交换 \(\nabla_\phi\) 与 \(\E_\varepsilon\) 就是交换求导与积分，本书第 12 部分\hyperref[sec:12-Real-Analysis-Support/04-Limits-and-Integration/03]{``求导穿过积分号需要控制差商''}给了完整条件，这里只做条件对账。

要用的结论是：若对几乎所有 \(\varepsilon\)，\(\phi\mapsto f(g_\phi(x,\varepsilon))\) 在 \(\phi\) 的一个邻域内可微，且差商被一个\emph{与 \(\phi\) 无关的可积函数}控制住，则换序成立。第一个条件排除的是 \(f\circ g\) 在某些 \(\varepsilon\) 上根本不可导（比如 \(f\) 有跳跃）；第二个条件排除的是差商虽然逐点收敛、却把质量漏到无穷远处——这正是控制收敛定理要挡的失败模式。VAE 的常规配置里两条都容易满足：\(g_\phi\) 是仿射的，decoder 是分段光滑的网络，高斯尾部提供了包络。

值得警惕的是另一类分解。取 \(z\sim \Uniform(0,\theta)\)，重参数化写得出来：\(z=\theta u\)，\(u\sim\Uniform(0,1)\)。可如果你不重参数化，直接对原式求导：

\[
\frac{d}{d\theta}\int_0^\theta \frac{1}{\theta}f(z)\,dz ,
\]

积分\emph{上限本身随参数移动}。此时``把 \(\nabla_\theta\) 挪进积分号''漏掉了边界项，答案是错的——正确的求导要用 Leibniz 法则补上 \(f(\theta)\) 那一项。支撑随参数移动的分布，就属于这种情形。重参数化在这里之所以还能用，恰恰是因为它把变动的支撑换成了固定的 \([0,1]\)，让参数只出现在被积函数里；而不是因为``换序总是对的''。这两件事必须分开记。

得分函数路线不受这个困扰：它对 \(q_\phi\) 求导时用的是同一个换序假设，但被积函数是 \(q_\phi f\)，只要 \(q_\phi\) 的支撑不随 \(\phi\) 变（大多数参数族满足），条件就是标准的。

\subsection{离散隐变量：分解不存在时怎么办}

重参数化的门槛现在看得很清楚：它\emph{要求存在}一个把参数与随机性分离的可微分解 \(z=g_\phi(\varepsilon)\)。高斯有，位置尺度族普遍有，Gamma 和 Dirichlet 要靠隐式微分或接受近似。类别分布没有——从 \(K\) 个类别里抽一个，输出是离散的，对概率参数的导数处处为零。

Gumbel-Softmax（也叫 concrete 分布）的做法是先给出一个精确的采样重写，再把它连续化。Gumbel-max 技巧说：若 \(g_1,\dots,g_K\) 独立取自标准 Gumbel 分布，则

\[
\argmax_k \big(\log \pi_k + g_k\big)
\]

的分布正是类别概率 \(\pi\) 的类别分布。这已经是一个重参数化了——随机性全在 \(g\) 里，参数在 \(\log\pi\) 里——只可惜 \(\argmax\) 不可导。于是把它换成带温度的 softmax：

\[
y_k = \frac{\exp\big((\log\pi_k+g_k)/\tau\big)}{\sum_{j}\exp\big((\log\pi_j+g_j)/\tau\big)}.
\]

\(y\) 是单纯形内部的一个连续点，对 \(\pi\) 和 \(\tau\) 都可微，反向传播畅通。

换来这一点的是它\emph{有偏}。你优化的不再是原来那个离散隐变量模型的目标，而是一个松弛版本的目标。温度 \(\tau\to 0\) 时 \(y\) 收敛到 one-hot，偏差消失；但同一个极限下 softmax 的雅可比趋于退化，梯度方差爆炸。这不是实现上的粗糙，是两端拉扯的真实权衡：偏差和方差被同一个旋钮反向控制，中间没有免费的位置。常见做法是把 \(\tau\) 从 1 附近退火到 0.5 左右后停住，或者用直通估计器在前向取 one-hot、反向用松弛的雅可比——后者偏差更大，但前向语义是干净的离散量。这些都是经验规律，不同任务上的最佳设定差别很大，没有理论保证。

选择因此可以归结成一句话：\(z\) 连续且分解存在，就用重参数化，方差优势是压倒性的；\(z\) 离散，先问自己能不能接受松弛带来的偏差，不能接受就退回得分函数估计并认真做基线。两条路都在算同一个梯度，只是一个从 \(f\) 的斜率取信息，一个从 \(f\) 的取值取信息。

\subsection{梯度通了，不等于隐变量有用}

这一节解决的是``梯度能不能到达 \(\phi\)''。它确实到达了：encoder 的参数开始更新，ELBO 开始下降，训练曲线看起来完全正常。

但把梯度接通与让隐变量真的承担信息，是两件独立的事。目标函数本身可能允许一个 encoder 输出与 \(x\) 无关的常数分布——梯度照样流过重参数化的那条边，只是它把 \(\mu_\phi\) 推向零、\(\sigma_\phi\) 推向一之后就停了。那时你有一条漂亮的损失曲线和一个什么都没编码的隐空间。下一节处理的正是这个退化解：它不是训练没跑好，而是 ELBO 的结构允许它存在。
